\section{Revisión de Literatura}
\label{sec:lit}

Esta sección resume seis trabajos representativos sobre reconstrucción de
batimetría y modelado de aguas someras con redes neuronales y métodos de
optimización, ordenados por relevancia para esta tesis. Cada uno se describe con
la misma estructura: problema, entrada y salida, datos, método, resultados clave
y relación con esta tesis. La Tabla~\ref{tab:lit-overview} los resume de un
vistazo.

%% IEEE-CUT: en la versión IEEE cada artículo se comprime a 1--2 frases; los bloques \paragraph{} detallados son solo para la tesis.

\begin{table}[H]
\centering
\caption{Panorama de los trabajos revisados.}
\label{tab:lit-overview}
\begin{tabular}{@{}llll@{}}
\toprule
Trabajo & Problema & Entrada $\to$ Salida & Método \\
\midrule
Dazzi (2024) & Directo               & BC, IC $\to (h,u,\zb)$      & PINN, MLP único \\
Ohara (2024) & Inverso conjunto      & $u,\eta \to (h,u,v,\zb,n)$  & PINN, red por campo \\
Tian (2025)  & Directo (formas)      & dominio $\to (h,u,v)$       & PINN, FNN única \\
Chen (2023)  & Inverso (costero)     & oleaje $\to \zb$            & PINN, ec.\ de oleaje \\
Angel (2024) & Inverso, datos reales & $\eta$ de sensores $\to \zb$ & Adjunto-EDP \\
Liu (2024)   & Inverso, dos etapas   & $\eta \to \zb$              & Surrogado CNN \\
\bottomrule
\end{tabular}
\end{table}

% ──────────────────────────────────────────────────────────────
\subsection{Dazzi et al.\ (2024) --- PINNs con SWE aumentado}
\label{sec:lit:dazzi}

\paragraph{Problema.}
\cite{dazzi2024} resuelven el problema \emph{directo} de las SWE con un sistema
\emph{aumentado} que incorpora la batimetría como variable: a las ecuaciones de
masa y momento se añaden dos relaciones. La primera, $\eta = h + \zb$, define la
elevación de la superficie libre como suma de profundidad y fondo, y acopla así
la incógnita ($\zb$) con la cantidad observable ($\eta$). La segunda,
$\pd{\zb}{t} = 0$, impone que el fondo es estacionario, lo que obliga a la red a
representar $\zb$ como un único campo espacial coherente con toda la evolución
temporal del flujo.

\paragraph{Entrada y salida.}
Recibe condiciones de borde e iniciales del flujo y produce los campos
$(h, u, \zb)$ de forma simultánea: la red aprende $\zb$ como parte de la
solución en lugar de recibirla como entrada.

\paragraph{Datos.}
Casos con solución analítica o de referencia: agua en reposo sobre fondo plano,
bump, escalón y parábola; flujos transitorios (perturbación, rotura de presa,
rotura de presa sobre escalón y oscilación de Thacker); y flujo estacionario
sobre un bump.

\paragraph{Método.}
Un único MLP (perceptrón multicapa) de 7 capas $\times$ 300 neuronas,
implementado en NVIDIA Modulus. La ecuación $\pd{\zb}{t}=0$ no se impone
estructuralmente, sino como un término más de la función de pérdida, con un peso
de penalización de 100 que fija cuán fuerte se castiga su incumplimiento frente
a los demás términos. La pendiente del fondo $\pd{\zb}{x}$ se obtiene por
diferenciación automática, consistente con el marco PINN.

\paragraph{Resultados clave.}
Reportan el error de cada variable por caso (su Tabla~2); para el caso del bump,
el RMSE (raíz del error cuadrático medio) de $\zb$ es $\approx 5{,}2$\,mm.
Entrenamiento de 30\,000 épocas de Adam; código disponible en Zenodo.
Limitación: dificultad con discontinuidades del fondo (escalón), donde la
predicción se suaviza o muestra picos locales, consistente con el sesgo
espectral de los PINNs.

\paragraph{Relación con esta tesis.}
Es la arquitectura de referencia: el diseño monolítico A2 del estudio
comparativo (Sección~\ref{sec:arquitecturas}) imita su red única, y la propuesta
de esta tesis (A1) la lleva a $\sim$18K parámetros mediante dos redes separadas,
imponiendo $\pd{\zb}{t}=0$ de forma estructural (la red de batimetría no recibe
$t$ como entrada) en lugar de como término de pérdida.

% ──────────────────────────────────────────────────────────────
\subsection{Ohara et al.\ (2024) --- Inversión conjunta de flujo y geometría}
\label{sec:lit:ohara}

\paragraph{Problema.}
\cite{ohara2024} emplean un PINN para invertir \emph{conjuntamente} el flujo y
la geometría del lecho de un cauce a partir de las SWE 2D de plano, recuperando
variables difíciles de medir sin etiquetas directas para ellas.

\paragraph{Entrada y salida.}
Como datos de entrenamiento recibe mediciones \emph{dispersas} de velocidad
superficial y de nivel de agua. Las redes, con entradas espaciales $(x,y)$,
producen los campos estacionarios de profundidad $h$, velocidades $u$ y $v$, la
cota del fondo $\zb$ y el coeficiente de rugosidad de Manning $n$.

\paragraph{Datos.}
La forma del fondo se midió en un experimento de laboratorio (canal con barras
de arena, mediante tomografía de flujo); el campo de flujo se generó con un
solver numérico 2D (Nays2D).

\paragraph{Método.}
Una red totalmente conectada \emph{separada por cada campo} ($h$, $u$, $v$,
$\zb$); el coeficiente de rugosidad $n$ se optimiza como un parámetro más. Usa
fricción de Manning e incorpora la conservación de masa (caudal constante por
sección transversal) como \emph{restricción redundante}, lo que mejora de forma
notable la convergencia y la precisión.

\paragraph{Resultados clave.}
Recupera la cota del fondo y la rugosidad sin datos de entrenamiento para esas
variables; el error de la cota del fondo es del orden del 7--12\,\% de la
profundidad de flujo uniforme. La precisión es robusta a observaciones
espacialmente dispersas siempre que su resolución supere la longitud de onda de
las formas del fondo.

\paragraph{Relación con esta tesis.}
Es el predecesor más directo: PINN, SWE e inversión conjunta de flujo y
geometría a partir de observaciones dispersas. Su diseño de \emph{una red
dedicada por variable} es la base del diseño A3 del estudio comparativo de
arquitecturas (Sección~\ref{sec:arquitecturas}). Esta tesis se diferencia en
que extiende el banco de pruebas al régimen transitorio 2D (Exp.~3 a 5),
compara sistemáticamente formas del residuo y arquitecturas de red, y evalúa
datos experimentales reales (Exp.~6 con Angel et al.\ \cite{angel2024}).

% ──────────────────────────────────────────────────────────────
\subsection{Tian et al.\ (2025) --- Formas de ecuación para PINNs-SWE}
\label{sec:lit:tian}

\paragraph{Problema.}
\cite{tian2025} comparan sistemáticamente formulaciones del residuo SWE para el
problema \emph{directo} en 2D --- con topografía y términos de lluvia ---, con
el fin de determinar qué forma de las ecuaciones produce mejor precisión y
convergencia en un PINN.

\paragraph{Entrada y salida.}
Recibe el dominio espacio-temporal $(x,y,t)$ y produce los campos de flujo
$(h,u,v)$. La topografía es un \emph{dato conocido}: funciones analíticas
explícitas cuyas derivadas se calculan analíticamente y entran en el término
fuente del residuo de la EDP.

\paragraph{Datos.}
Once casos de prueba: estáticos, con lluvia y dinámicos (rotura de presa 2D y
circular, problema mareal y mareal con lluvia), con soluciones analíticas o de
volúmenes finitos como referencia.

\paragraph{Método.}
El modelo es una \emph{única} red totalmente conectada (5 capas $\times$ 300
neuronas, activación Tanh) con activación adaptativa por neurona (N-LAAF). Se
comparan tres formas del residuo: primitiva $(h,u,v)$; primitiva-conservativa
$\bm{R} = \bm{A}\,\bm{r}$ con
$\bm{A} = \begin{psmallmatrix} 1 & 0 & 0 \\ u & h & 0 \\ v & 0 & h \end{psmallmatrix}$,
que pondera automáticamente los residuos sin hiperparámetros adicionales; y
conservativa $(h,hu,hv)$. El manejo seco-mojado usa una función de Heaviside.

\paragraph{Resultados clave.}
La forma \emph{primitiva-conservativa} resulta la más fiable; la forma de la
ecuación afecta de manera significativa la precisión del PINN. Incorporar la
condición de entropía (conservación de energía) no mejora la precisión y a veces
desestabiliza el entrenamiento. Los errores van de $\sim$$10^{-4}$\,m en los
casos estáticos a $\sim$$10^{-2}$\,m en rotura de presa. Al revisar su código
publicado, esta tesis identificó un error de implementación en la forma
primitiva (un factor $h$ adicional en el gradiente de presión) que el artículo
no menciona.

\paragraph{Relación con esta tesis.}
Fundamenta el estudio de ablación de formas del residuo
(Sección~\ref{sec:ablation}). A diferencia de esta tesis, Tian resuelve el
problema \emph{directo} --- la topografía es un dato conocido, no una incógnita
---, por lo que no forma parte del estudio comparativo de arquitecturas.

% ──────────────────────────────────────────────────────────────
\subsection{Chen et al.\ (2023) --- PINN para batimetría costera}
\label{sec:lit:chen}

\paragraph{Problema.}
\cite{chen2023} aplican un PINN al problema \emph{inverso} de mapear la
batimetría somera costera y reconstruir el campo de oleaje de forma simultánea,
a partir de observaciones dispersas de las olas. Es, según los autores, la
primera vez que se aplican PINNs a la inversión de batimetría en aguas someras.

\paragraph{Entrada y salida.}
Recibe observaciones dispersas del oleaje --- número de onda y/o altura de ola
y, en una de sus variantes, mediciones escasas de profundidad --- y produce la
batimetría (profundidad de agua) junto con el campo de oleaje completo.

\paragraph{Datos.}
Playas con barras (perfiles idealizados) con datos sintéticos del modelo XBeach,
y un bajo circular con datos \emph{experimentales} de laboratorio (Chawla
et al., 1996). El ensayo con datos reales de campo se deja como trabajo futuro.

\paragraph{Método.}
PINN cuya restricción física \emph{no} son las SWE, sino la ecuación de balance
de energía del oleaje junto con la relación de dispersión (lineal o no lineal).
Usa redes totalmente conectadas compuestas (una para el campo de energía del
oleaje, otra para el número de onda y la profundidad), de 4 capas $\times$ 30
neuronas, activación Tanh, entrenadas con Adam y L-BFGS-B.

\paragraph{Resultados clave.}
Recupera la batimetría con alta precisión --- RMSE de la profundidad del orden
de $10^{-2}$\,m y $R^2 > 0{,}98$ en los casos de prueba --- y supera con claridad
a una red neuronal convencional sin la restricción física. Emplear la relación
de dispersión \emph{no lineal} es decisivo: con la versión lineal el error de la
profundidad se degrada en la zona de rompientes.

\paragraph{Relación con esta tesis.}
Es un PINN de inversión de batimetría en el dominio costero, afín a la
motivación de monitoreo de esta tesis. Su física (ecuaciones de oleaje) difiere
de las SWE, por lo que sirve como antecedente del paradigma pero no como
arquitectura comparable en el estudio de la Sección~\ref{sec:arquitecturas}.

% ──────────────────────────────────────────────────────────────
\subsection{Angel et al.\ (2024) --- Inversión con datos experimentales}
\label{sec:lit:angel}

\paragraph{Problema.}
\cite{angel2024} presentan la primera reconstrucción de batimetría a partir de
datos \emph{experimentales} reales, mediante optimización restringida por EDP
(ecuaciones en derivadas parciales), no con PINNs.

\paragraph{Entrada y salida.}
Recibe las series temporales de elevación de la superficie libre registradas por
sensores puntuales y produce el perfil de batimetría $\zb(x)$.

\paragraph{Datos.}
Canal de oleaje de 12\,m de longitud, con un bump de 200\,mm en agua de 0,3\,m,
4 sensores y 20 réplicas experimentales. Código y datos abiertos.

\paragraph{Método.}
Adjunto continuo de las SWE 1D para el cálculo eficiente del gradiente; SWE
discretizadas con Dedalus (espectral, $M = 68$ puntos de Chebyshev); descenso de
gradiente con line search de Armijo.

\paragraph{Resultados clave.}
NRMSE (RMSE normalizado) de $\approx$14--16\,\% reconstruyendo desde los
sensores; usar 2 sensores da una reconstrucción casi tan precisa como con 3. La
posición del máximo de la batimetría se recupera bien, aunque su altura se
subestima.

\paragraph{Relación con esta tesis.}
Es el \textbf{baseline clásico de restricción dura} contra el que se compara el
PINN (Exp.\ 6, Sección~\ref{sec:comp}); su éxito allí donde el PINN de
penalización blanda falla motiva el diagnóstico blando-vs-duro de este trabajo.

% ──────────────────────────────────────────────────────────────
\subsection{Liu et al.\ (2024) --- Surrogado CNN para inversión}
\label{sec:lit:liu}

\paragraph{Problema.}
\cite{liu2024} adoptan un enfoque en dos etapas: primero entrenan una red que
\emph{emula} las SWE 2D (un sustituto rápido y diferenciable del solver
numérico) y luego recuperan la batimetría optimizándola contra ese sustituto.

\paragraph{Entrada y salida.}
A partir de observaciones de elevación de la superficie del agua produce la
batimetría $\zb$; el sustituto intermedio emula el mapa directo, de $\zb$ a los
campos de flujo.

\paragraph{Datos.}
Unas 3000 simulaciones del solver SRH-2D sobre un canal sintético rectangular
(25,6\,m $\times$ 6,4\,m) para entrenar el sustituto.

\paragraph{Método.}
El sustituto es un autoencoder convolucional (CNN, red neuronal convolucional)
entrenado para predecir los campos de flujo a partir de la batimetría. Una vez
entrenado, sus pesos se \emph{congelan} y la batimetría pasa a ser la variable a
optimizar: partiendo de una estimación inicial, se ajusta $\zb$ por descenso de
gradiente (Adam) para minimizar el desajuste entre la superficie predicha por el
sustituto y la observada, más regularización de suavidad y de rango. Los
gradientes respecto de $\zb$ se obtienen por diferenciación automática
retropropagando a través del sustituto ya entrenado, sin reentrenarlo.

\paragraph{Resultados clave.}
RMSE de $\zb$ de $\approx$0,09\,m a partir solo de observaciones de elevación de
superficie (con regularización; el peor caso, sin regularización, llega a
$\approx$0,21\,m).

\paragraph{Relación con esta tesis.}
Paradigma alternativo (surrogado diferenciable frente a PINN): rápido en la
inversión, pero inaplicable sin un gran conjunto de simulaciones previas, a
diferencia del enfoque PINN, que no requiere datos de simulación.
